\documentclass[12pt,a4paper,twocolumn]{article} 

% Packages
\usepackage[utf8]{inputenc}
\usepackage[T1]{fontenc}
\usepackage{graphicx}
\usepackage{amsmath,amssymb}
\usepackage{booktabs}
\usepackage{array}
\usepackage{cite}
\usepackage{url}
\usepackage{hyperref}
\usepackage{float}
\usepackage{geometry}

\geometry{margin=1in}

% Title Information
\title{Building a Khmer Traffic Sign Dataset for Automated Detection and Classification}
\author{
    Contributor 1 \and
    Contributor 2 \and
    Contributor 3 \and
    Contributor 4
}
\date{\today}

\begin{document}

\maketitle

% Abstract
\begin{abstract}
Real-time data warehousing has become increasingly important for organizations
that require up-to-date information to support business decision-making.
Traditional batch ETL approaches introduce latency and limit the ability of
organizations to respond quickly to operational changes. This paper proposes
a real-time data warehouse architecture using Change Data Capture (CDC),
Apache Kafka, Apache Spark, and PostgreSQL.
\end{abstract}

\tableofcontents
\newpage

% Sections
\include{sections/introduction}
\include{sections/literature_review}
\section{Methodology}
The following section discusses the techniques utilized to develop and create the Khmer traffic sign detection and classification approach.
The overall research process and experimental procedures are describe in the following section.
\subsection{Methodology Workflow}
The proposed methodlogy establishes a systematic workflow for Khmer traffic sign detection and classification.
The process begins with data preparation and data preprocessing for the annotated traffic sign dataset, followed by the model development
and training under experimental set up. After training, the models will evaluated on the test dataset using standard object detection metrics.
The obtained results are subsequently compared to identify model that achieves performance for Khmer dataset. 
\subsection{Data Preprocessing}

\subsection{Model Development}
\subsection{Model Training}
\subsection{Evaluation}
\include{sections/system_architecture}
\include{sections/implementation}
\include{sections/results}
\section{Conclusion}
This study explore the application of YOLO-based object detection for detected and classified Khmer traffic sign.
The experimental results on the Khmer dataset demonstrated both YOLOv8 and YOLOv10 models achieved strong detection
performance, with YOLOv8 showing a small overall advantage across the evaluation measure. YOLOv8 was able to detect
and classify Khmer traffic signs even when the signs had different sizes and the images had different lighting conditions
and image quality.

This study contributes to the development of Khmer traffic sign recognition by evaluating modern object detection models
using a Khmer traffic sign dataset. The results demonstrate the potential of deep learning for automatically detecting and
classifying Khmer traffic signs. Furthermore, this study provides a useful reference for future research and the development
of computer vision-based traffic sign recognition systems for the Cambodian road environment.

For future work, the Khmer traffic sign dataset can be expanded by including more traffic signs images within each classes 
and incorporating additional traffic sign classes to create more comprehensive dataset. Future research could also explore other advanced and lightweight
object detection models to enhance detection performance in real-world conditions including weather conditions, lighting environments,
and road locations. 

% References
\bibliographystyle{IEEEtran}
\bibliography{references/references}

\end{document}
